\documentclass[10pt,a4paper]{amsart}
\usepackage[margin=19mm]{geometry}
\usepackage{amsmath,amssymb,mathtools}
\usepackage[scr=rsfs,cal=euler]{mathalfa}
\usepackage{xfrac}
\usepackage[unicode,hidelinks,bookmarks=false]{hyperref}
\newcommand{\ie}{i.e.\ }
\newcommand{\cf}{cf.\ }
\newcommand{\resp}{respectively\ }
\newcommand{\Subsection}{\subsection}
\providecommand{\nobreakdash}{\nobreak\hbox{-}}
\setlength{\emergencystretch}{2em}
\multlinegap0pt
\usepackage{bm}
\usepackage{bbm}
\usepackage{dsfont}
\usepackage{xspace}
\usepackage{cancel}
\usepackage{mathtools}
\usepackage{amsthm}
\makeatletter
\@ifundefined{theorem}{\newtheorem{theorem}{Theorem}}{}
\@ifundefined{Fact}{\newtheorem{Fact}[theorem]{Observation}}{}
\makeatother
\title[On the representation number of grid graphs and cylindric gr]{On the representation number of grid graphs and cylindric grid graphs}
\author{Nawaf Shafi Alshammari, Sergey Kitaev, Artem Pyatkin}
\date{Source: arXiv:2507.16469v1 (CC BY 4.0)}
\begin{document}
\maketitle

\noindent\emph{Source and licence.} On the representation number of grid graphs and cylindric grid graphs. Version arXiv:2507.16469v1 (CC BY 4.0), distributed under the Creative Commons Attribution 4.0 International licence, as recorded in the arXiv licence metadata for this exact version and shown on the abstract page https://arxiv.org/abs/2507.16469v1. Full rights notice: licenses/arxiv-2507.16469-CC-BY-4.0.txt.\par\smallskip

\noindent\textit{Editorial scope.} Counted unit: the Theorem whose source label is \texttt{cyl4} in the article (source0.tex lines 330--330, source label \texttt{cyl4}), delivered together with the article's own complete original proof of it (source0.tex lines 332--398, one \texttt{proof} environment of 67 lines). No mathematics is written, repaired or re-derived: statement and proof are the article's own text line for line. Required context retained uncounted: evident (lines 180--182). Every definition, display and label of those lines is retained unchanged. Omitted from this package and not redistributed: 1--179, 183--329, 331--331, 399--556. Nothing inside a retained excerpt is omitted or altered: every retained body is delivered exactly as the article's lines are written, so no omission receipt of any kind accompanies this package. Citations: the retained excerpts carry no citation command at all, so no bibliography entry of the article is transcribed and no reference is redistributed. Reference adaptation: none was needed and none was made; every reference inside the delivered unit resolves inside it, so no unresolved reference or citation is produced. Printed slips kept literal: no printed word, symbol, hypothesis or number of the counted statement or of its proof was altered. Layout and class: the article's own class is \texttt{article} with packages including amsmath, amssymb, amsfonts, amsthm, tikz, graphicx, placeins, tkz-graph; the wrapper is the lane's pinned \texttt{amsart} editorial wrapper at 10pt on A4, which restates the theorem-like environments the retained text needs and adds the article's own macro definitions that the retained text needs (none), with extra wrapper packages (bm, bbm, dsfont, xspace, cancel, mathtools). The delivered numbering is the unit's own. Assets: no retained span contains a figure environment, a graphics-inclusion command, a PDF-page inclusion, a table or a TikZ picture; no image file is copied into this package, the package declares no asset dependency and no image file is redistributed with it. Licence: version arXiv:2507.16469v1 of this article is distributed under the Creative Commons Attribution 4.0 International licence, as recorded in the arXiv licence metadata for this exact version and shown on the abstract page https://arxiv.org/abs/2507.16469v1. A full rights notice is delivered at licenses/arxiv-2507.16469-CC-BY-4.0.txt. Characters: every retained excerpt is pure ASCII, so the delivered engine, XeLaTeX with its default Latin Modern text font, reports no missing character.
\begin{Fact}~\label{evident}
If $xy\in E$ then $x^i>y^{i-1}$ and $y^i>x^{i-1}$ for all $i$.
\end{Fact}

\begin{theorem}\label{cyl4}  Cylindrical grid graph $\mbox{CGr}_{m,n}$ is $3$-representable for all $m\geq 1$ and $n\ge 4$.\end{theorem}

\begin{proof}
For each $m$ we construct a 3-uniform word $w_m$ reprsenting $\mbox{CGr}_{m,n}$ and let $W_n=w_n\setminus w_{n-1}$. In our construction $W_m$ coincides, after replacing $x_{m,i}$ with  $x_i$ for all $j$, with \texttt{Od} if $m$ is odd and with \texttt{Ev} if $m$ is even. Namely, we start with
$$w_1=x_{11}^1x_{1n}^1x_{12}^1x_{11}^2\ldots x_{1,n-1}^1x_{1,n-2}^2x_{1,n}^2x_{1,n-1}^2x_{11}^3\ldots x_{1,n-2}^3x_{1,n}^3x_{1,n-1}^3$$
and apply one of the following constructions for $w_m$ depending on the parity of $m\ge 2$. 

If $m$ is even, we substitute:
\begin{eqnarray*}
x_{m-1,j}^2 & \mbox{ by } & x_{m,j}^1x_{m-1,j}^2  \mbox{ for all } j=1,\ldots,n;\\
x_{m-1,1}^3 & \mbox{ by } & x_{m,1}^2x_{m,n}^2x_{m-1,1}^3; \\
x_{m-1,j}^3 & \mbox{ by } & x_{m,j}^2x_{m,j-1}^3x_{m-1,j}^3 \mbox{ for all } j=2,\ldots,n-2;\\
x_{m-1,n}^3 & \mbox{ by } & x_{m,n-1}^2x_{m,n-2}^3x_{m-1,n}^3;\\
x_{m-1,n-1}^3 & \mbox{ by } & x_{m,n}^3x_{m-1,n-1}^3x_{m,n-1}^3.
\end{eqnarray*}

If $m$ is odd, we substitute:
\begin{eqnarray*}
x_{m-1,1}^2 & \mbox{ by } & x_{m,1}^1x_{m-1,1}^2;\\
x_{m-1,n}^2 & \mbox{ by } & x_{m,n}^1x_{m-1,n}^2 ;\\
x_{m-1,j}^2 & \mbox{ by } & x_{m,j}^1x_{m,j-1}^2x_{m-1,j}^2 \mbox{ for all } j=2,\ldots,n-1;\\
x_{m-1,n}^3 & \mbox{ by } & x_{m,n}^2x_{m,n-1}^2x_{m,1}^3x_{m-1,n}^3;\\
x_{m-1,n-1}^3 & \mbox{ by } & x_{m,2}^3x_{m,3}^3\ldots x_{m,n-2}^3x_{m,n}^3x_{m-1,n-1}^3x_{m,n-1}^3.
\end{eqnarray*}

Let us show that $w_m$ indeed represents $\mathrm{CGr}_{m,n}$. Clearly, $W_m$ coincides with \texttt{Od} and \texttt{Ev} for odd and even values of $m$, respectively; hence, the vertices $x_{m,1}, \ldots, x_{m,n}$ induce the cycle $C_n$. We have to verify that $x_{m,j}$ is adjacent in $w_{m-1}$ to  $x_{m-1,j}$ only for all $ j=1,\ldots,n$. 
Throughout the proof we will use the following easy observation: if $x<y$ and there is a factor $zy$ then $x<z$.
Consider two cases. \\[-3mm]

\noindent
{\bf Case 1.} Let $m\ge 3$ be odd. Note that $w_m$ contains a factor  $ x_{m,1}^1x_{m-1,1}^2$ while $w_{m-1}$ contains a factor $x_{m-1,1}^2x_{m-1,n}^2x_{m-2,1}^3$ (remind that  $w_{m-1}$ was constructed by the rules of even $m-1$). Therefore, $[x_{m,1}^1,x_{m-2,1}^3]\cap w_{m-2}=\emptyset.$ By property~3), $x_{m-2,1}^3>x_{m-2,j}^2$, and thus,
 $x_{m,1}^1>x_{m-2,j}^2$ for all $j=1,\ldots,n$. Hence, by Observation~\ref{evident}, no $x_{m,j}$ is adjacent to any $x\in w_{m-2}$.

Clearly, $[x_{m,n}^2,x_{m,n}^3]\cap w_{m-1}=\{x_{m-1,n}^3\}$; so, $x_{m,n}$ maybe adjacent in $w_{m-1}$ only to $x_{m-1,n}$. We have $[x_{m,1}^1,x_{m,1}^2]\cap w_{m-1} =
\{ x_{m-1,1}^2,x_{m-1,n}^2 \}$, but $x_{m,1}^3<x_{m-1,n}^3$. So, $x_{m,1}$ and $x_{m-1,n}$ are not adjacent. Similarly, 
$[x_{m,n-1}^2,x_{m,n-1}^3]\cap w_{m-1} = \{ x_{m-1,n}^3,x_{m-1,n-1}^3 \}$, but $x_{m-1,n}^2<x_{m,n-1}^1$ and hence  $x_{m,n-1}$ and $x_{m-1,n}$ are not adjacent. 
Let $j\in [2,n-2]$. Note that in  \texttt{Ev} we have $[x_j^2,x_{j+1}^2] =\{ x_{j-1}^3\}$. Therefore, $[x_{m,j}^1,x_{m,j}^2]\cap w_{m-1} =\{ x_{m-1,j}^2,x_{m-1,j-1}^3\}$.
% Since for every $j=2,\ldots, n-2$  in  \texttt{Ev} $[x_j^2,x_{j+1}^2] =\{ x_{j-1}^3\}$, we have $[x_{m-1,j}^1,x_{m,j}^2]\cap w_{m-1} =\{ x_{m-1,j}^2,x_{m-1,j-1}^3\}$.
%there is a factor $x_{m,j}^1x_{m,j-1}^2x_{m-1,j}^2$ in $w_m$; in particular, the  inequality $x_{m,j}^2<x_{m-1,j+1}^2$ holds for all $j=1,\ldots, n-2$. However, in $W_{m-1}$ (that coincides with Ev) we have$[x_{m-1,j}^2,x_{m-1,j+1}^2] =\{ x_{m-1,j-1}^3\}$ and  
Since by Observation~\ref{evident}, $x_{m,j}$ cannot be adjacent in to $x_{m-1,j-1}$, the only possible neighbor of  $x_{m,j}$ in $w_{m-1}$ is $x_{m-1,j}$.

It remains to verify that the edges $x_{m,j}x_{m-1,j}$ exist for all $j=1,\ldots,n$. Clearly,  $x_{m-1,j}^3<x_{m,j}^3$ for all $j$. Since
$[x_{m,j}^1,x_{m-1,j}^2]\cap w_{m-1}$ is either empty or contains $x_{m,j-1}^2$ only, we have $x_{m-1,j}^1<x_{m,j}^1<x_{m-1,j}^2<x_{m,j}^2$. 
The inequalities $x_{m,n}^2<x_{m-1,n}^3$ and  $x_{m,n-1}^2<x_{m-1,n-1}^3$ are evident. Finally, for all $j=1,\ldots,n-2$ by property~(2) we have
$x_{m,j}^2<x_{m-1,j+1}^2<x_{m-1,j}^3$.  So, $x_{m,j}x_{m-1,j}\in E$  for all $j=1,\ldots, n$. \\[-3mm]

\noindent
{\bf Case 2. } Let $m\ge 2$ be even. If $m\ge 4$ then by the previous case, $x_{m,1}^1>x_{m-1,1}^1>x_{m-3,j}^2$ for all $j$ and so, there are no edges between 
 $x_{m,i}$ and $w_{m-3}$. 

Note that $w_m$ contains a factor  $x_{m,1}^2x_{m,n}^2x_{m-1,1}^3$ while $w_{m-1}$ contains a factor $x_{m-1,1}^3x_{m-2,n}^3$. Hence,
 for all $i=1,\ldots,n$ and $j=1,\ldots,n-2$ we have $x_{m,i}^2>x_{m,1}^2>x_{m-2,j}^3$ and by Observation~\ref{evident}, $x_{m,i}x_{m-2,j}\not\in E$. 
Since $w_m$ has a factor $x_{m,1}^1x_{m-1,1}^2$ and $w_{m-1}$ has factors  $x_{m-1,1}^2x_{m-2,2}^2$ and $x_{m-1,n}^1x_{m-2,n}^2$, and $x_n^1<x_2^2$ in \texttt{Od}, we derive
$x_{m,i}^1>x_{m,1}^1>x_{m-2,n}^2$ and  by Observation~\ref{evident}, $x_{m,i}x_{m-2,n}\not\in E$ for all $i=1,\ldots,n$. By property~3) we have
$x_{m-2,n-1}^1<x_{m-2,1}^2<x_{m-1,1}^2$. Since there is a factor $x_{m,1}^1x_{m-1,1}^2$ in $w_m$, for all $i=1,\ldots,n$ the inequality $x_{m-2,n-1}^1<x_{m,i}^1$ holds.
Since $w_{m-1}$ ends with a factor $x_{m-2,n-1}^3x_{m-1,n-1}^3$ and $x_{m-1,n-1}^3$ was substituted by $ x_{m,n}^3x_{m-1,n-1}^3x_{m,n-1}^3$ in $w_m$, we have $x_{m-2,n-1}^3>x_{m,i}^3$ for all $i=1,\ldots, n-2$ and  by definition, $x_{m,i}x_{m-2,n-1}\not\in E$ for these $i$. Finally, $w_{m-1}$ and $w_m$ have factors 
$ x_{m-1,n-1}^1x_{m-1,n-2}^2x_{m-2,n-1}^2$ and $x_{m,n}^1x_{m-1,n}^2$, respectively. By property~(2), $x_{m-1,n-2}^1<x_{m-1,n}^2$ and hence by property~(1)
$x_{m-2,n-1}^2<x_{m,n}^1<x_{m,n-1}^1$. So, by Observation~\ref{evident}, $x_{m-2,n-1}$ cannot be adjacent to $x_{m,n}$ and $x_{m,n-1}$. So, no edges connect $W_m$ with $w_{m-2}$.

Since \texttt{Od} contains the permutation $x_1^3x_2^3\ldots x_{n-2}^3x_n^3x_{n-1}^3$, we have $[x_{m,j}^2,x_{m,j}^3]\cap W_{m-1}=\{x_{m-1,j}^3\}$ for all $j=1,\ldots,n-2$. 
Clearly,  $[x_{m,n}^1,x_{m,n}^2]\cap W_{m-1}=\{x_{m-1,n}^2,x_{m-1,n-1}^2\}$, but $x_{m,n}^3<x_{m-1,n-1}^3$, and thus $x_{m,n}x_{m-1,n-1}\not\in E$.
Finally,  $[x_{m,n-1}^2,x_{m,n-1}^3]\cap W_{m-1}=\{x_{m-1,n}^3,x_{m-1,n-1}^3\}$. By property~(1), $x_{m-1,n}^2<x_{m-1,n-1}^2$. Since $w_m$ has a factor
%$x_{m,n}^1,x_{m-1,n}^2$ and 
$x_{m,n-1}^1x_{m-1,n-1}^2$, we have $x_{m-1,n}^2<x_{m,n-1}^1$. So,  by Observation~\ref{evident}, $x_{m,n-1}x_{m-1,n}\not\in E$.

It remains to verify that the edges $x_{m,j}x_{m-1,j}$ exist for all $j=1,\ldots,n$. The factors $x_{m,j}^1x_{m-1,j}^2$ imply that 
$x_{m-1,j}^1<x_{m,j}^1<x_{m-1,j}^2<x_{m,j}^2$ for all $j$. The inequalities $x_{m,j}^2<x_{m-1,j}^3<x_{m,j}^3$ can be seen directly from the construction. So, 
 $x_{m,j}x_{m-1,j}\in E$ for all $j=1,\ldots, n$.
\end{proof}

\end{document}
